The gardener's maintenance cycle is the recurring editorial routine that keeps the manuscript aligned as the surrounding work changes. Rather than waiting for a build failure or a human complaint, the gardener periodically inspects the current state of the repository and compares each affected section against the rest of the book. The goal is not merely to detect errors, but to identify the smallest useful editorial action that will restore coherence.

A maintenance pass begins with change inspection. The gardener reviews newly edited sections, updated outlines, revised prompts, and any downstream artifacts that may have been invalidated by those changes. This first pass is intentionally broad: a local edit can alter terminology, dependencies, or claims in sections that do not themselves contain the edit. The gardener therefore treats the manuscript as a connected system rather than a set of isolated files.

After identifying the changed region, the gardener compares the section against several reference points. The parent section provides the higher-level argument and establishes what the subsection is supposed to contribute. Sibling sections provide immediate narrative context and help reveal duplication, gaps, or inconsistent terminology. Declared dependencies indicate which earlier material the section relies on, while downstream uses show where the section's concepts, definitions, or phrasing are reused later in the book. These comparisons let the gardener distinguish a purely local wording issue from a structural mismatch that will propagate elsewhere.

This comparison step is what makes the maintenance cycle editorial rather than mechanical. A section may be internally well formed yet still repeat a claim already made in a sibling, omit a transition that the parent section needs, or introduce a term that later sections assume without definition. The gardener looks for those mismatches and classifies them by scope. Some issues can be repaired by a single section agent. Others require coordinated revisions across multiple sections. Still others are not prose problems at all, but missing support that must be routed to the references agent.

Once the gardener has identified the relevant issue class, it queues targeted work instead of rewriting the manuscript wholesale. A repeated claim may become a revision request to the section that introduced it. A thin transition may be assigned to the section that needs a bridging sentence. A missing definition may be routed to the section where the term first appears, or to the references agent if the term requires external support. This targeted routing keeps the maintenance cycle small, inspectable, and reversible.

The recurring nature of the cycle matters as much as the checks themselves. Because the manuscript evolves through many small agent actions, coherence cannot be enforced only at the end of drafting. The gardener therefore acts as a standing editorial process: inspect, compare, classify, and queue. In effect, it keeps the book in a state of continuous maintenance, where local edits are evaluated against the broader structure before they harden into global drift.

This section's role in the larger chapter is to establish that gardening is not a passive validation step. It is an active control loop that watches the manuscript's moving parts, compares them across structural boundaries, and dispatches the smallest corrective task that will preserve the book's internal consistency. The next section separates those local revision requests into concrete categories, showing how repeated claims, thin transitions, missing definitions, and unsupported assertions each route to a different repair path.
